\chapter{Shaders and the .fx Bytecode Gap}
\label{ch:shader-fx-gap}

Chapter~\ref{ch:what-is-cna} named this the single biggest real gap in CNA. This short
chapter explains precisely what the gap is, why it is architectural rather than an
oversight, and what does exist today as a real, working alternative.

\section{What real XNA actually ships}

A real XNA game never ships \texttt{.fx} shader source text. The XNA content pipeline
compiles \texttt{.fx} effect files to a proprietary compiled bytecode format at build time,
and the compiled bytecode --- not the HLSL source --- is what ships inside the game's
\texttt{.xnb} content files and gets loaded at runtime via
\texttt{new Effect(GraphicsDevice, byte[])}. On the real Xbox 360/Windows XNA runtime, that
bytecode is interpreted directly; FNA instead ships MojoShader, a from-scratch bytecode
\emph{parser and cross-compiler} that turns the same compiled Direct3D~9 shader bytecode into
GLSL (or other target shading languages) at load time, which is exactly what lets FNA run
real, unmodified, historical XNA content unchanged.

\section{Why \texttt{Effect(GraphicsDevice\&, byte[])} always throws in CNA}

CNA has no equivalent of MojoShader. Its \cnaclass{Effect} constructor that takes a raw
bytecode buffer --- the one a real \texttt{.xnb}-loaded custom effect would call --- always
throws \texttt{NotImplementedException} (Chapter~\ref{ch:stock-effects}). This is not a
missing feature that simply has not been gotten to yet in the way, say, a missing
\texttt{SurfaceFormat} is; it is a genuinely large, separate compiler-engineering project ---
parsing Direct3D~9 shader bytecode's own instruction format and re-targeting it to
GLSL/SPIR-V/WGSL/HLSL-for-D3D11-12 for each of CNA's own backends is comparable in scope to
building a second MojoShader from scratch. The project's own planning documents
(\texttt{docs/shader-effect-vs-fx-bytecode.md} and \texttt{docs/fx-bytecode-support-plan.md})
track this as a distinct, deliberately-scoped future phase rather than a task sitting in an
ordinary backlog.

\section{The measurable cost of the gap}

This gap has a precise, checkable cost, not a vague one: it is the reason 23 of the 86
official XNA samples ported in the sibling \texttt{cna-samples} repository
(Chapter~\ref{ch:ecosystem-map}) do not run on CNA today. Any sample whose content pipeline
compiles a custom \texttt{.fx} effect --- rather than relying only on the five stock effects
from Chapter~\ref{ch:stock-effects} --- hits this exact constructor and throws. This is also
precisely the number this book's migration-guide chapter uses to help you judge, before
starting a port, whether an existing XNA/FNA game is likely to need this gap closed or not:
a game built entirely on stock effects and \cnaclass{SpriteBatch} ports cleanly; a game with
even one custom \texttt{.fx} shader compiled into its content does not, yet.

\section{What does exist today: ShaderEffect, a runtime-compiled alternative}

CNA does not leave custom shading entirely unaddressed --- it takes a different path than
bytecode compatibility. A \texttt{ShaderEffect} class exists on the three backends where it
has been built out (\textbf{D3D9}, \textbf{D3D11}, \textbf{D3D12}, all covered in
Chapter~\ref{ch:direct3d-backends}), letting a game supply HLSL \emph{source} at runtime,
compiled on the spot via \texttt{D3DCompile()} (or, on \texttt{D3D9} specifically, the same
real \texttt{d3dcompiler\_47.dll} the stock-effect pipeline itself uses, isolated into its
own link target so the rest of the \texttt{D3D9} backend can stay free of a
\texttt{d3dcompiler} dependency). This is not the same thing as bytecode compatibility ---
a game must supply new HLSL source, not reuse an existing compiled \texttt{.fx} blob --- but
it is a real, working, pixel-verified path for adding custom shading to a CNA game today,
provided you are targeting one of the three Direct3D backends and are willing to author (or
port) the shader as source rather than load a precompiled asset.

\cnaclass{GpuDrawParams} (Chapter~\ref{ch:stock-effects}) carries a
\texttt{customEffectBackend} pointer specifically to route a \texttt{ShaderEffect}-compiled
program through the same per-draw dispatch every stock effect uses, and the backend contract
layer's \cnaclass{IEffectBackend} interface (Chapter~\ref{ch:backend-architecture}) exposes
\texttt{CompileProgram(vertSrc, fragSrc)} plus a full set of default-no-op uniform setters ---
the same interface a stock effect's compiled shader program is bound through, reused rather
than duplicated for the custom-shader case.

\section{The honest summary}

If your porting target relies only on \cnaclass{SpriteBatch} and the five stock effects from
Chapter~\ref{ch:stock-effects}, this gap does not touch you at all. If it relies on a custom
compiled \texttt{.fx} shader loaded from content, it does --- today, on every backend except
the three where \texttt{ShaderEffect} exists, and even there only if you can supply the
shader as HLSL source rather than a precompiled binary. This is precisely the kind of
qualified, specific claim Chapter~\ref{ch:what-is-cna} promised this book would preserve:
not ``shaders don't work,'' and not ``shaders are fully supported,'' but a named, bounded gap
with a measured cost (23 of 86 samples) and a real, if partial, alternative already built.
